% Options for packages loaded elsewhere
\PassOptionsToPackage{unicode}{hyperref}
\PassOptionsToPackage{hyphens}{url}
\documentclass[
  a4paper,11pt]{article}
\usepackage{xcolor}
\usepackage[margin=1in]{geometry}
\usepackage{amsmath,amssymb}
\usepackage{graphicx}
\setcounter{secnumdepth}{-\maxdimen} % remove section numbering
\usepackage{iftex}
\ifPDFTeX
  \usepackage[T1]{fontenc}
  \usepackage[utf8]{inputenc}
  \usepackage{textcomp}
\else
  \usepackage{unicode-math}
  \defaultfontfeatures{Scale=MatchLowercase}
  \defaultfontfeatures[\rmfamily]{Ligatures=TeX,Scale=1}
\fi
\usepackage{lmodern}
\ifPDFTeX\else
  \setmainfont[]{Times New Roman}
\fi
\IfFileExists{upquote.sty}{\usepackage{upquote}}{}
\IfFileExists{microtype.sty}{%
  \usepackage[]{microtype}
  \UseMicrotypeSet[protrusion]{basicmath}
}{}
\makeatletter
\@ifundefined{KOMAClassName}{%
  \IfFileExists{parskip.sty}{%
    \usepackage{parskip}
  }{%
    \setlength{\parindent}{0pt}
    \setlength{\parskip}{6pt plus 2pt minus 1pt}}
}{% if KOMA class
  \KOMAoptions{parskip=half}}
\makeatother
\usepackage{longtable,booktabs,array}
\newcounter{none}
\usepackage{calc}
\usepackage{etoolbox}
\makeatletter
\patchcmd\longtable{\par}{\if@noskipsec\mbox{}\fi\par}{}{}
\makeatother
\IfFileExists{footnotehyper.sty}{\usepackage{footnotehyper}}{\usepackage{footnote}}
\makesavenoteenv{longtable}
\usepackage{bookmark}
\IfFileExists{xurl.sty}{\usepackage{xurl}}{}
\urlstyle{same}
\usepackage{hyperref}
\hypersetup{
  hidelinks,
  pdfcreator={LaTeX via pandoc}}
\setlength{\emergencystretch}{3em}
\providecommand{\tightlist}{%
  \setlength{\itemsep}{0pt}\setlength{\parskip}{0pt}}
\usepackage{amsthm}
\newtheorem{theorem}{Theorem}[section]
\newtheorem{proposition}[theorem]{Proposition}
\newtheorem{lemma}[theorem]{Lemma}
\newtheorem{corollary}[theorem]{Corollary}
\theoremstyle{definition}
\newtheorem{definition}[theorem]{Definition}
\theoremstyle{remark}
\newtheorem{remark}[theorem]{Remark}
\newtheorem{observation}[theorem]{Observation}
\begin{document}
\section{Interactions of Shape-Invariant Waves: Axis-Directional
Displacement Transfer, Wave Packet Deformation, and Retroactive
Construction of
Causality}\label{interactions-of-shape-invariant-waves-axis-directional-displacement-transfer-wave-packet-deformation-and-retroactive-construction-of-causality}

\textbf{Author}: Noriaki Kihara\\
\textbf{Affiliation}: WF System Co., Ltd.\\
\textbf{ORCID}: 0009-0004-6753-4020\\
\textbf{Version}: v1.0\\
\textbf{Date}: April 2026\\
\textbf{DOI}: 10.5281/zenodo.19709800

\begin{center}\rule{0.5\linewidth}{0.5pt}\end{center}

\subsection{Abstract}\label{abstract}

Building on the three-mode classification of shape-invariant waves
introduced in the companion paper {[}3{]} (direction-constrained waves
\(\kappa = 2\), spherical waves \(\kappa = 0\), standing waves
\(\kappa = 1\)), we provide a unified description of interactions
between different modes. The fundamental mechanism of interaction is
defined as \textbf{axis-directional displacement transfer} -- the change
of a direction-constrained wave's state vector
\((\mathbf{k}, \Delta\mathbf{k})\) along the direction of the non-zero
axis carried by a spherical wave -- and the four forces
(electromagnetism, weak force, strong force, gravity) are unified as
applications of this single operation to different axes. Coupling with a
standing wave (\(\kappa = 1\)) modifies the propagation speed ratio
\(|k_x / k_t|\) of a direction-constrained wave, providing a geometric
origin for mass. Under the conserved quantity \(\int |\psi|^2 \, dV\),
energy transfer deforms the amplitude and width of a wave packet
inversely, and this deformation gives a dynamical description of ``wave
packet collapse.'' Interactions mediated by spherical waves with
\(k_t = 0\) carry no intrinsic causal ordering, and the shortest
geodesic path and the localization of the wave packet are retroactively
constructed from the interaction outcome.

\textbf{Keywords}: shape-invariant wave, axis-directional displacement
transfer, wave packet deformation, origin of mass, retroactive
construction of causality, conservation law

\begin{center}\rule{0.5\linewidth}{0.5pt}\end{center}

\subsection{S1 Introduction and Scope of
Claims}\label{s1-introduction-and-scope-of-claims}

This paper provides a unified description of interactions between waves
of different modes within the framework of shape-invariant waves
constructed in the companion papers {[}1{]}--{[}3{]}, and as
consequences derives the origin of mass, the classification of forces,
wave packet deformation, and causal structure geometrically.

The claims of this paper are limited to the following five points:

\begin{enumerate}
\def\labelenumi{\arabic{enumi}.}
\tightlist
\item
  All interactions are described uniformly as displacement transfer
  along the non-zero axes of spherical waves. The difference among types
  of force reduces solely to the difference in the axis along which the
  transfer occurs.
\item
  Coupling with a standing wave (\(\kappa = 1\)) modifies the speed
  ratio of a direction-constrained wave, providing a geometric origin
  for mass.
\item
  Under the conserved quantity \(\int |\psi|^2 \, dV\), energy transfer
  deterministically deforms the shape of a wave packet, and this
  deformation corresponds to ``wave packet collapse'' in observation.
\item
  Interactions mediated by spherical waves with \(k_t = 0\) carry no
  intrinsic causal ordering, and the path and localization are
  retroactively constructed from the outcome.
\item
  The sign of the interaction (attraction / repulsion) between
  direction-constrained waves is determined geometrically by the
  relative sign of their \(t\)-axis components.
\end{enumerate}

The following elements are \textbf{outside the scope of this paper}:

\begin{itemize}
\tightlist
\item
  Numerical derivation of coupling constants (specific values of
  \(\alpha\), \(G_F\), \(\alpha_s\), \(G\))
\item
  Quantitative computation of scattering cross sections and decay widths
\item
  Renormalization and scale dependence of running coupling constants
\item
  Rigorous reconstruction of gauge groups in the continuum limit
\end{itemize}

\textbf{Note on physical nomenclature}: This paper uses physical names
(photon, graviton, electron, flavor, color charge, etc.) assuming
correspondence with Table A of {[}4{]}, but this is for interpretive
convenience; the geometric content of each proposition does not depend
on the physical identification. All propositions are derived solely from
the component structure of wave vectors (\(\kappa\), \(k_t\), type of
axis).

\begin{center}\rule{0.5\linewidth}{0.5pt}\end{center}

\subsection{S2 Unified Definition of
Interaction}\label{s2-unified-definition-of-interaction}

\subsubsection{S2.1 Two-Layer Structure}\label{s2.1-two-layer-structure}

In the tangent bundle \(\mathbb{Z}^4 \times \mathbb{Z}^4\) introduced in
{[}2{]} S8, the state of a direction-constrained wave (\(\kappa = 2\))
is:

\[(\mathbf{k},\; \Delta\mathbf{k}) \in \mathbb{Z}^4 \times \mathbb{Z}^4\]

\begin{itemize}
\tightlist
\item
  \(\mathbf{k}\): \textbf{static label} (type of wave, corresponding to
  Table A of {[}4{]})
\item
  \(\Delta\mathbf{k}\): \textbf{dynamic state} (corresponding to energy
  and momentum)
\end{itemize}

\subsubsection{S2.2 Axis-Directional Displacement
Transfer}\label{s2.2-axis-directional-displacement-transfer}

\textbf{Definition 2.1} (Axis-directional displacement transfer): When a
spherical wave (\(\kappa = 0\)) interacts with a direction-constrained
wave (\(\kappa = 2\)), the state of the direction-constrained wave
changes as follows:

\[(\mathbf{k},\; \Delta\mathbf{k}) \;\longrightarrow\; (\mathbf{k} + \delta\mathbf{k},\; \Delta\mathbf{k} + \delta(\Delta\mathbf{k}))\]

where the displacements \(\delta\mathbf{k}\) and
\(\delta(\Delta\mathbf{k})\) satisfy the following conditions:

\textbf{(I1) Axis-directionality}: The non-zero components of
\(\delta\mathbf{k}\) and \(\delta(\Delta\mathbf{k})\) exist only in the
direction of the non-zero axis of the spherical wave in Table A of
{[}4{]}.

\textbf{(I2) Conservation law}: The total sum of the tangent bundle
states of all waves involved in the interaction is conserved:

\[\sum_i (\mathbf{k}_i + \Delta\mathbf{k}_i) = \text{const}\]

\textbf{(I3) Wave packet conservation}: The conserved quantity
\(\int |\psi|^2 \, dV\) of each wave is individually conserved.

\subsubsection{S2.3 Static Label Change and Dynamic State
Change}\label{s2.3-static-label-change-and-dynamic-state-change}

Axis-directional displacement transfer encompasses two qualitatively
distinct types of change:

{\def\LTcaptype{none} % do not increment counter
\begin{longtable}[]{@{}
  >{\raggedright\arraybackslash}p{(\linewidth - 6\tabcolsep) * \real{0.3438}}
  >{\raggedright\arraybackslash}p{(\linewidth - 6\tabcolsep) * \real{0.1875}}
  >{\raggedright\arraybackslash}p{(\linewidth - 6\tabcolsep) * \real{0.3125}}
  >{\raggedright\arraybackslash}p{(\linewidth - 6\tabcolsep) * \real{0.1562}}@{}}
\toprule\noalign{}
\begin{minipage}[b]{\linewidth}\raggedright
Type of change
\end{minipage} & \begin{minipage}[b]{\linewidth}\raggedright
Target
\end{minipage} & \begin{minipage}[b]{\linewidth}\raggedright
Physical meaning
\end{minipage} & \begin{minipage}[b]{\linewidth}\raggedright
Example
\end{minipage} \\
\midrule\noalign{}
\endhead
\bottomrule\noalign{}
\endlastfoot
\(\delta\mathbf{k} \neq 0\) & Static label \(\mathbf{k}\) & Conversion
of wave type & Flavor transition, color charge exchange \\
\(\delta(\Delta\mathbf{k}) \neq 0\) & Dynamic state \(\Delta\mathbf{k}\)
& Change in energy/momentum & Scattering, acceleration \\
\end{longtable}
}

The type of spherical wave determines which (or both) of these is
dominant.

\subsubsection{S2.4 Unification of the Four
Forces}\label{s2.4-unification-of-the-four-forces}

From the non-zero axis structure of spherical waves in Table A of
{[}4{]}, the four forces are described as applications of the same
operation to different axes:

{\def\LTcaptype{none} % do not increment counter
\begin{longtable}[]{@{}
  >{\raggedright\arraybackslash}p{(\linewidth - 6\tabcolsep) * \real{0.1333}}
  >{\raggedright\arraybackslash}p{(\linewidth - 6\tabcolsep) * \real{0.2333}}
  >{\raggedright\arraybackslash}p{(\linewidth - 6\tabcolsep) * \real{0.3000}}
  >{\raggedright\arraybackslash}p{(\linewidth - 6\tabcolsep) * \real{0.3333}}@{}}
\toprule\noalign{}
\begin{minipage}[b]{\linewidth}\raggedright
Force
\end{minipage} & \begin{minipage}[b]{\linewidth}\raggedright
Spherical wave
\end{minipage} & \begin{minipage}[b]{\linewidth}\raggedright
Non-zero axis
\end{minipage} & \begin{minipage}[b]{\linewidth}\raggedright
Dominant change
\end{minipage} \\
\midrule\noalign{}
\endhead
\bottomrule\noalign{}
\endlastfoot
Electromagnetism & \(\gamma\; (0,0,0,0,R,0)\) & \(R\) &
\(\delta(\Delta\mathbf{k})\) (momentum change) \\
Weak force & \(W^{\pm}\; (0,0,0,\pm t,0,0)\) & \(t\) &
\(\delta\mathbf{k}\) (label change) \\
Strong force & \(g\; (0,0,0,0,0,Q_i)\) \({}^{*}\) & \(Q\) &
\(\delta\mathbf{k}\) (label change) \\
Gravity & \(G\; (0,0,0,\pm t,\pm R,0)\) & \(t, R\) &
\(\delta(\Delta\mathbf{k})\) (energy change) \\
\end{longtable}
}

\({}^{*}\) While \(\gamma\), \(W^{\pm}\), and \(G\) are spherical waves
with static non-zero components on the \(t\) or \(R\) axis, \(g\) is a
spherical wave carrying a discrete label value \(Q_i\) on the \(Q\) axis
that mediates color charge exchange by transitioning to a wave with a
different label value \(Q_j\). That is, the ``non-zero axis'' of \(g\)
acts not as a fixed value but as a change in the \(Q\) label (detailed
in S6).

\textbf{Proposition 2.1}: All four forces are described as
axis-directional displacement transfer per Definition 2.1, and the
difference among types of force reduces solely to the difference in the
axis along which the transfer occurs.

\begin{center}\rule{0.5\linewidth}{0.5pt}\end{center}

\subsection{S3 Origin of Mass: Coupling with Standing
Waves}\label{s3-origin-of-mass-coupling-with-standing-waves}

\subsubsection{S3.1 Geometric Structure of
Coupling}\label{s3.1-geometric-structure-of-coupling}

The standing wave mode (\(\kappa = 1\), {[}3{]} S5) fills the entire
subjective space \(S^3\) as a fundamental mode along one spatial axis. A
direction-constrained wave (\(\kappa = 2\)) propagates through this
standing wave background.

A direction-constrained wave \(\mathbf{k} = (k_x, 0, 0, k_t)\) and a
standing wave \(\mathbf{k}_H = (k_{H}, 0, 0, 0)\) share the spatial
\(x\)-axis. On this shared axis, the phase structures of the waves
overlap, giving rise to coupling.

\subsubsection{S3.2 Modification of the Speed
Ratio}\label{s3.2-modification-of-the-speed-ratio}

The propagation speed of a direction-constrained wave without coupling
is \(v = k_x / k_t\). Coupling with a standing wave affects the phase
structure in the \(x\)-axis direction, modifying the effective spatial
component:

\[k_x \;\longrightarrow\; k_x^{\text{eff}} = k_x - \delta k_H\]

where \(\delta k_H > 0\) is proportional to the coupling strength with
the standing wave. The sign is negative (subtraction) because the
background phase structure of the standing wave overlaps in antiphase
with the spatial phase of the direction-constrained wave, coupling in
the direction that cancels the effective spatial phase variation. The
propagation speed after coupling is:

\[v^{\text{eff}} = \frac{k_x^{\text{eff}}}{k_t} = \frac{k_x - \delta k_H}{k_t} < \frac{k_x}{k_t} = v\]

The speed decreases -- this is the geometric origin of \textbf{mass}.

\textbf{Proposition 3.1}: Coupling with a standing wave (\(\kappa = 1\))
reduces the effective spatial component of a direction-constrained wave,
restricting its propagation speed to \(v < c\). This speed restriction
is observed as mass.

\subsubsection{S3.3 Waves That Do Not
Couple}\label{s3.3-waves-that-do-not-couple}

Spherical waves (\(\kappa_s = 0\)) have no spatial components. Although
a standing wave has a phase structure along one spatial axis, there is
no shared spatial axis with a spherical wave.

\textbf{Proposition 3.2}: Spherical waves with \(\kappa_s = 0\) have no
coupling axis with a standing wave, yielding \(\delta k_H = 0\). Their
speed remains \(v = c\) without modification.

This is consistent with the fact that gauge bosons with \(\kappa_s = 0\)
(\(\gamma\), \(g\)) are massless in Table A of {[}4{]}. Although
\(W^{\pm}\) and \(Z^0\) have \(\kappa_s = 0\) and thus lack the direct
coupling via spatial axis described in Proposition 3.2, they can acquire
mass through an indirect standing wave coupling mechanism via the \(t\)
or \(R\) axis. Proposition 3.2 addresses only coupling through direct
sharing of a spatial axis; the details of indirect coupling mechanisms
via the \(t\) or \(R\) axis are outside the scope of this paper.

\subsubsection{S3.4 Mass Hierarchy}\label{s3.4-mass-hierarchy}

The coupling strength \(\delta k_H\) depends on the overlap integral
between the direction-constrained wave and the standing wave on the
spatial axis. Due to the generation structure (selection of spatial axes
\(k_1, k_2, k_3\)) shown in {[}3{]} S3.3, the overlap integrals can
differ among the three generations.

On \(S^3\), the three spatial axes are locally equivalent under
\(\mathrm{SO}(3)\) symmetry, but they can become non-equivalent in their
global coupling with the standing wave mode. This suggests a geometric
origin for mass differences between generations, but the derivation of
specific mass ratios is outside the scope of this paper.

\begin{center}\rule{0.5\linewidth}{0.5pt}\end{center}

\subsection{S4 Electromagnetic Interaction and
Scattering}\label{s4-electromagnetic-interaction-and-scattering}

\subsubsection{S4.1 Displacement Transfer by
Photons}\label{s4.1-displacement-transfer-by-photons}

The photon \(\gamma = (0, 0, 0, 0, +, 0)\) is a spherical wave with a
non-zero component on the \(R\) axis. The \(R\) axis corresponds to the
radial direction of central projection, and on \(\mathbb{Z}^4\) it acts
as a change in the dynamic state \(\Delta\mathbf{k}\).

When a direction-constrained wave (electron type) \((k_x, 0, 0, k_t)\)
absorbs a photon:

\[\Delta\mathbf{k} \;\longrightarrow\; \Delta\mathbf{k} + \delta(\Delta\mathbf{k})_{\gamma}\]

The static label \(\mathbf{k}\) does not change (the electron remains an
electron). Only the dynamic state changes:

\begin{itemize}
\tightlist
\item
  Change of momentum direction (scattering angle)
\item
  Change of energy (absorption of photon energy)
\end{itemize}

\textbf{Proposition 4.1}: Interaction with a photon preserves the static
label of a direction-constrained wave and changes only the dynamic state
\(\Delta\mathbf{k}\).

\subsubsection{S4.2 Electromagnetic Interaction Between
Direction-Constrained
Waves}\label{s4.2-electromagnetic-interaction-between-direction-constrained-waves}

The electromagnetic interaction between two direction-constrained waves
\(\psi_1\), \(\psi_2\) is described as photon exchange:

\[\psi_1 \;\xrightarrow{\text{photon emission}}\; \psi_1' + \gamma \;\xrightarrow{\text{photon absorption}}\; \psi_1' + \psi_2'\]

By conservation law (I2):

\[\Delta\mathbf{k}_1 + \Delta\mathbf{k}_2 = \Delta\mathbf{k}_1' + \Delta\mathbf{k}_2'\]

\subsubsection{S4.3 Sign of Interaction: Attraction and
Repulsion}\label{s4.3-sign-of-interaction-attraction-and-repulsion}

The relative sign of the \(t\)-axis components of two
direction-constrained waves determines the directionality of the
interaction.

The \(t\)-axis component \(k_t\) corresponds to the charge structure in
Table A of {[}4{]}. Photon exchange between two waves with same-sign
\(k_t\) acts repulsively, while opposite-sign \(k_t\) acts attractively.

{\def\LTcaptype{none} % do not increment counter
\begin{longtable}[]{@{}llll@{}}
\toprule\noalign{}
\(k_t\) of \(\psi_1\) & \(k_t\) of \(\psi_2\) & Sign product &
Interaction \\
\midrule\noalign{}
\endhead
\bottomrule\noalign{}
\endlastfoot
\(+\) & \(+\) & \(+1\) & Repulsion \\
\(+\) & \(-\) & \(-1\) & Attraction \\
\(-\) & \(-\) & \(+1\) & Repulsion \\
\end{longtable}
}

\textbf{Proposition 4.2}: The sign of the electromagnetic interaction
between two direction-constrained waves is determined geometrically by
the relative sign of their \(t\)-axis components,
\(\mathrm{sign}(k_{t,1} \cdot k_{t,2})\). The specific correspondence
``same sign = repulsion, opposite sign = attraction'' is an
interpretation layer dependent on the correspondence between \(k_t\) and
charge sign in Table A of {[}4{]}; the geometric content of the
proposition is that the relative sign uniquely determines the
directionality of the interaction.

\subsubsection{S4.4 Pair Annihilation}\label{s4.4-pair-annihilation}

A direction-constrained wave and its antiparticle have all components
with opposite signs (\(C\) transformation in {[}3{]} S6.5):

\[\psi: (k_x, 0, 0, k_t) \quad \longleftrightarrow \quad \bar{\psi}: (-k_x, 0, 0, -k_t)\]

When both interact in the same lattice region, the sum of the static
labels vanishes:

\[(k_x, 0, 0, k_t) + (-k_x, 0, 0, -k_t) = (0, 0, 0, 0)\]

A pair of \(\kappa = 2\) waves transitions to a \(\kappa = 0\) state. By
conservation law (I2), the energy of the vanished static labels is
emitted as spherical waves (photons).

\textbf{Proposition 4.3}: Pair annihilation of a direction-constrained
wave and an anti-direction-constrained wave is the process in which a
pair of \(\kappa = 2\) waves is converted into \(\kappa = 0\) spherical
waves, made possible by the vanishing of the sum of static labels.

\begin{center}\rule{0.5\linewidth}{0.5pt}\end{center}

\subsection{S5 Gravitational Interaction and Wave Packet
Deformation}\label{s5-gravitational-interaction-and-wave-packet-deformation}

\subsubsection{S5.1 Displacement Transfer by
Gravitons}\label{s5.1-displacement-transfer-by-gravitons}

The graviton \(G = (0, 0, 0, \pm t, \pm R, 0)\) has non-zero components
on both the \(t\) axis and the \(R\) axis. By condition (I1),
displacement transfer acts in both the \(t\) and \(R\) directions:

\[\delta(\Delta\mathbf{k})_G = (\delta(\Delta k)_t,\; \delta(\Delta k)_R)\]

This means simultaneous change of energy (\(t\) component) and spatial
configuration (\(R\) component).

\subsubsection{S5.2 Universal Coupling}\label{s5.2-universal-coupling}

We recall the description of acceleration introduced in {[}3{]} S8. The
second-order difference of the child-space central projection determines
the acceleration:

\[\Delta^2\mathbf{k}(\mathbf{k}_t) = D^2\pi(\mathbf{k}_t) \cdot \Delta\mathbf{k}_t\]

Graviton exchange is the discrete realization of this
\(\Delta^2\mathbf{k}\). Since gravitons act on the dynamic state
\(\Delta\mathbf{k}\):

\begin{itemize}
\tightlist
\item
  They do not depend on the static label \(\mathbf{k}\) (regardless of
  the type of wave)
\item
  They act on all waves with \(\Delta\mathbf{k} \neq 0\) (all waves
  carrying energy)
\end{itemize}

\textbf{Proposition 5.1}: Interaction with gravitons acts only on the
dynamic state \(\Delta\mathbf{k}\) and does not act on the static label
\(\mathbf{k}\). Therefore all shape-invariant waves with
\(\Delta\mathbf{k} \neq 0\) couple to gravitons (universal coupling).

\subsubsection{S5.3 Wave Packet
Deformation}\label{s5.3-wave-packet-deformation}

Energy transfer through interaction with gravitons deforms the shape of
a wave packet. Under the conserved quantity
\(\int |\psi|^2 \, dV = E_0\) ({[}1{]} S5, {[}2{]} S5, condition (C1)),
the amplitude \(A\) and characteristic width \(w\) of the wave packet
satisfy:

\[A^2 \cdot w^d = E_0 = \text{const} \quad (d = 3: \text{spatial dimension of } S^3)\]

For a wave packet that receives energy:

\[\Delta E > 0 \;\Rightarrow\; A \uparrow,\; w \downarrow \quad (\text{amplitude increase, width decrease})\]

\[\Delta E < 0 \;\Rightarrow\; A \downarrow,\; w \uparrow \quad (\text{amplitude decrease, width increase})\]

\textbf{Proposition 5.2}: Under the conserved quantity
\(\int |\psi|^2 \, dV\), energy transfer deforms the amplitude and width
of a wave packet inversely. An increase in energy is accompanied by
localization (decrease in width) of the wave packet.

\subsubsection{S5.4 Wave Packet Deformation of Spherical
Waves}\label{s5.4-wave-packet-deformation-of-spherical-waves}

Spherical waves (\(\kappa = 0\)) also couple to gravitons (Proposition
5.1: they couple if \(\Delta\mathbf{k} \neq 0\)). When a spherical wave
receives energy:

\begin{itemize}
\tightlist
\item
  The amplitude of the spherical wave increases and its spatial extent
  decreases
\item
  In the extreme limit, the spherical wave becomes point-like and
  localized
\end{itemize}

This corresponds to a geometric description of the gravitational lensing
effect. When a spherical wave (photon) passes through the gravitational
field of a massive body, the wave packet is deformed by interaction with
gravitons, causing the propagation path to bend.

\begin{center}\rule{0.5\linewidth}{0.5pt}\end{center}

\subsection{S6 Strong Interaction and
Confinement}\label{s6-strong-interaction-and-confinement}

\subsubsection{\texorpdfstring{S6.1 \(Q\)-Axis
Transition}{S6.1 Q-Axis Transition}}\label{s6.1-q-axis-transition}

The gluon \(g = (0, 0, 0, 0, 0, Q_i \to Q_j)\) is defined as a
transition on the \(Q\) axis. By condition (I1), displacement transfer
acts only in the \(Q\)-axis direction:

\[\delta\mathbf{k}_g: \quad Q \;\longrightarrow\; Q'\]

The spatial components, time component, and \(R\) component of the
direction-constrained wave do not change. Only the discrete label \(Q\)
changes.

\subsubsection{S6.2 Three-Bit Structure of Color
Charge}\label{s6.2-three-bit-structure-of-color-charge}

In Table A of {[}4{]}, \(Q\) is encoded as three bits
\((c_1, c_2, c_3)\):

\begin{itemize}
\tightlist
\item
  \(c_2 c_3\): color charge (\(00\) = colorless, \(01, 10, 11\) = three
  colors)
\item
  \(c_1\): weak isospin
\end{itemize}

The \(Q\) transition of a gluon is a change in the \(c_2 c_3\) bits,
corresponding to color charge exchange. Transitions among the three
non-zero values (\(01, 10, 11\)) of \(c_2 c_3\) give
\(3 \times 3 - 1 = 8\) independent gluons.

\subsubsection{S6.3 Confinement
Condition}\label{s6.3-confinement-condition}

In the framework of child-space central projection introduced in {[}3{]}
S8, forming an independent subjective space requires that the collection
of waves be \textbf{color-neutral}.

\textbf{Definition 6.1} (Color neutrality condition): A collection of
direction-constrained waves \(\{\psi_i\}\) is color-neutral if the
composition of the \(c_2 c_3\) bits reduces to \((0, 0)\). Here the
composition of bits is defined as bitwise exclusive OR (XOR, addition
over \(\mathbb{F}_2\)):

\[(a_1, a_2) \oplus (b_1, b_2) = (a_1 \oplus b_1,\; a_2 \oplus b_2) \quad (\oplus: \text{mod}\; 2 \text{ addition})\]

Realizations of color neutrality:

{\def\LTcaptype{none} % do not increment counter
\begin{longtable}[]{@{}lll@{}}
\toprule\noalign{}
Configuration & Color combination & Name \\
\midrule\noalign{}
\endhead
\bottomrule\noalign{}
\endlastfoot
Three waves & \((01) \oplus (10) \oplus (11) = (00)\) & Baryon type \\
Two waves & \((c_2 c_3) \oplus (c_2 c_3) = (00)\) & Meson type \\
One wave & \((c_2 c_3) = (00)\) & Lepton type \\
\end{longtable}
}

\textbf{Proposition 6.1}: A collection of direction-constrained waves
that does not satisfy the color neutrality condition cannot form an
independent subjective space and is confined within a color-neutral
collection.

\subsubsection{S6.4 Asymptotic Freedom}\label{s6.4-asymptotic-freedom}

Inside a color-neutral collection (inside the shared subjective space),
the curvature of the child-space central projection is small, and
direction-constrained waves move nearly freely. As the boundary of the
collection is approached, the projection curvature increases, making it
impossible to cross the boundary.

\textbf{Proposition 6.2}: The interior of a color-neutral collection has
small projection curvature (asymptotic freedom), while at the boundary
the projection curvature diverges (confinement).

\begin{center}\rule{0.5\linewidth}{0.5pt}\end{center}

\subsection{S7 Weak Interaction and Flavor
Transition}\label{s7-weak-interaction-and-flavor-transition}

\subsubsection{\texorpdfstring{S7.1 \(t\)-Axis
Displacement}{S7.1 t-Axis Displacement}}\label{s7.1-t-axis-displacement}

\(W^{\pm} = (0, 0, 0, \pm t, 0, 0)\) is a spherical wave with a non-zero
component on the \(t\) axis. By condition (I1):

\[\delta\mathbf{k}_W: \quad k_t \;\longrightarrow\; k_t' = k_t + \delta k_t\]

The change in the \(t\)-axis component corresponds to the up-type /
down-type conversion in Table A of {[}4{]}. Simultaneously, the \(c_1\)
bit of the \(Q\) axis also changes.

\subsubsection{S7.2 Concrete Example of Flavor
Transition}\label{s7.2-concrete-example-of-flavor-transition}

Transformation of a direction-constrained wave by \(W^+\) absorption:

\[\underbrace{(k_x, 0, 0, -k_t, 0, Q_d)}_{\text{down-type}} + \underbrace{(0, 0, 0, +, 0, 0)}_{W^+} \;\longrightarrow\; \underbrace{(k_x, 0, 0, +k_t, 0, Q_u)}_{\text{up-type}}\]

The sign of the \(t\) axis is inverted (\(-k_t \to +k_t\)), and the
isospin bit \(c_1\) of \(Q\) changes.

\subsubsection{\texorpdfstring{S7.3 \(k_t \neq 0\) and CP
Violation}{S7.3 k\_t \textbackslash neq 0 and CP Violation}}\label{s7.3-k_t-neq-0-and-cp-violation}

As shown in Proposition 6.4 of {[}3{]}, CP symmetry can be violated only
in interactions with \(k_t \neq 0\). \(W^{\pm}\) has
\(k_t = \pm 1 \neq 0\), and is the only one among the four types of
spherical waves to have \(k_t \neq 0\).

\textbf{Proposition 7.1}: Only \(W^{\pm}\), which has a non-zero
component on the \(t\) axis, mediates CP symmetry violation. This is a
direct consequence of \(k_t \neq 0\) introducing asymmetry in the time
direction.

\subsubsection{S7.4 Chirality Selection
Rule}\label{s7.4-chirality-selection-rule}

For the chirality \(\chi = \mathrm{sign}(k_{s_0} \cdot k_t)\)
(\(k_{s_0}\): non-zero spatial component) defined in Definition 6.1 of
{[}3{]}, the \(t\)-axis structure of \(W^{\pm}\) imposes a coupling
condition with a specific chirality.

\(W^{\pm} = (0, 0, 0, \pm t, 0, 0)\) has a phase structure only on the
\(t\) axis. Coupling with a direction-constrained wave
\((k_{s_0}, 0, 0, k_t)\) is established only when the \(t\)-axis phase
of \(W^{\pm}\) and the \(t\)-axis phase of the direction-constrained
wave have opposite signs
(\(\mathrm{sign}(k_t) \neq \mathrm{sign}(k_{t,W})\)). By Definition 6.1
of {[}3{]}, \(\chi = \mathrm{sign}(k_{s_0} \cdot k_t)\), so this
coupling condition corresponds to \(\chi = -1\) (left-handed).

\textbf{Proposition 7.2}: Coupling with \(W^{\pm}\) is restricted to
direction-constrained waves with chirality \(\chi = -1\) (left-handed).
Right-handed (\(\chi = +1\)) waves do not couple with \(W^{\pm}\).

\begin{center}\rule{0.5\linewidth}{0.5pt}\end{center}

\subsection{S8 Wave Packet Deformation and Retroactive Construction of
Causality}\label{s8-wave-packet-deformation-and-retroactive-construction-of-causality}

\subsubsection{S8.1 Conservation-Law Description of Wave Packet
Deformation}\label{s8.1-conservation-law-description-of-wave-packet-deformation}

We generalize the mechanism of wave packet deformation introduced in
S5.3. In any interaction where a direction-constrained wave receives
energy, under the conserved quantity \(\int |\psi|^2 \, dV = E_0\):

\[A_{\text{after}}^2 \cdot w_{\text{after}}^d = A_{\text{before}}^2 \cdot w_{\text{before}}^d = E_0\]

Concentration of energy density (increase in \(A\)) necessarily
accompanies spatial localization (decrease in \(w\)).

\subsubsection{S8.2 Dynamical Description of
``Observation''}\label{s8.2-dynamical-description-of-observation}

In conventional quantum mechanics, ``wave packet collapse upon
observation'' is introduced axiomatically. In the present framework:

\begin{enumerate}
\def\labelenumi{\arabic{enumi}.}
\tightlist
\item
  A direction-constrained wave (\(\kappa = 2\)) interacts with a
  spherical wave (\(\kappa = 0\))
\item
  Energy transfer increases the amplitude of the wave packet
\item
  By conservation law (I3), the width decreases (localization)
\item
  A sufficiently localized wave packet is described as ``detected as a
  particle''
\end{enumerate}

``Observation'' is not a special operation but merely an instance of
wave packet deformation due to interaction. No external ``observer'' is
required, and the deformation of the wave packet is a deterministic
process governed by conservation laws.

\textbf{Proposition 8.1}: Wave packet collapse is a
conservation-law-governed wave packet deformation due to interaction,
and there is no need to postulate a special process called
``observation.''

\subsubsection{S8.3 Retroactive Construction of
Causality}\label{s8.3-retroactive-construction-of-causality}

As shown in Proposition 4.2 of {[}3{]}, spherical waves with \(k_t = 0\)
are time-reversal symmetric and carry no intrinsic causal ordering. The
``deterministic'' in Proposition 8.1 and the ``retroactive'' below do
not contradict each other. They belong to different layers:

\begin{itemize}
\tightlist
\item
  \textbf{Deterministic}: Conservation laws (I2)(I3) hold strictly, and
  the outcome of wave packet deformation is uniquely determined
\item
  \textbf{Retroactive}: Because \(k_t = 0\) entails no causal ordering
  (no temporal order of ``cause then effect''), which event is the
  ``cause'' and which is the ``effect'' remains indeterminate
\end{itemize}

That is, the structure of ``strictly obeying conservation laws, yet
possessing no causal ordering'' arises naturally from the time-reversal
symmetry of \(k_t = 0\). Combining this property with wave packet
deformation:

\textbf{(R1) Retroactive determination of path}: A spherical wave
(\(k_t = 0\)) propagates isotropically in all directions (Proposition
4.1 of {[}3{]}). Viewed from the outcome of an interaction, the
spherical wave \textbf{appears to have} propagated along the shortest
geodesic. Since no causal ordering exists, the outcome retroactively
determines the path.

\textbf{(R2) Retroactive determination of localization}: Wave packet
deformation is a deterministic process governed by conservation laws
(Proposition 8.1). However, because \(k_t = 0\) entails no causal
ordering, ``the interaction occurred, therefore localization happened''
and ``localization occurred, therefore the interaction appears to have
happened'' are indistinguishable.

\textbf{(R3) Principle of least action}: As the synthesis of (R1) and
(R2), the variational principle (principle of least action) of natural
law is derived not as an axiom but as a consequence of the causal
symmetry of interactions mediated by spherical waves with \(k_t = 0\).

\textbf{Proposition 8.2}: In interactions mediated by spherical waves
with \(k_t = 0\), (a) the shortest path of propagation, (b) the
localization of the wave packet, and (c) the variational principle are
all retroactively constructed from the absence of causal ordering.

\subsubsection{\texorpdfstring{S8.4 The Case of
\(k_t \neq 0\)}{S8.4 The Case of k\_t \textbackslash neq 0}}\label{s8.4-the-case-of-k_t-neq-0}

Interactions mediated by \(W^{\pm}\) (\(k_t \neq 0\)) are not causally
symmetric. (R1)--(R3) do not apply:

\begin{itemize}
\tightlist
\item
  The path is causally determined (emission first, absorption later)
\item
  Time-reversal symmetry is broken
\item
  CP non-conservation becomes possible (Proposition 7.1)
\end{itemize}

{\def\LTcaptype{none} % do not increment counter
\begin{longtable}[]{@{}
  >{\raggedright\arraybackslash}p{(\linewidth - 8\tabcolsep) * \real{0.3125}}
  >{\raggedright\arraybackslash}p{(\linewidth - 8\tabcolsep) * \real{0.1875}}
  >{\raggedright\arraybackslash}p{(\linewidth - 8\tabcolsep) * \real{0.1875}}
  >{\raggedright\arraybackslash}p{(\linewidth - 8\tabcolsep) * \real{0.1042}}
  >{\raggedright\arraybackslash}p{(\linewidth - 8\tabcolsep) * \real{0.2083}}@{}}
\toprule\noalign{}
\begin{minipage}[b]{\linewidth}\raggedright
\(k_t\) of spherical wave
\end{minipage} & \begin{minipage}[b]{\linewidth}\raggedright
Causal ordering
\end{minipage} & \begin{minipage}[b]{\linewidth}\raggedright
Path selection
\end{minipage} & \begin{minipage}[b]{\linewidth}\raggedright
CP
\end{minipage} & \begin{minipage}[b]{\linewidth}\raggedright
Corresponding force
\end{minipage} \\
\midrule\noalign{}
\endhead
\bottomrule\noalign{}
\endlastfoot
\(k_t = 0\) & None & Retroactive (geodesic) & Conserved &
Electromagnetism, strong force \\
\(k_t \neq 0\) & Present & Causal & Can be violated & Weak force \\
\end{longtable}
}

\begin{center}\rule{0.5\linewidth}{0.5pt}\end{center}

\subsection{S9 Summary}\label{s9-summary}

We summarize the main results of this paper:

\begin{enumerate}
\def\labelenumi{\arabic{enumi}.}
\item
  \textbf{Unified interaction mechanism}: All forces are described as
  displacement transfer along the non-zero axes of spherical waves via
  (I1)--(I3), and the type of force reduces solely to the difference in
  the axis (S2).
\item
  \textbf{Origin of mass}: Coupling with a standing wave
  (\(\kappa = 1\)) decreases the speed ratio \(|k_x / k_t|\) of a
  direction-constrained wave, producing mass as \(v < c\). Spherical
  waves with \(\kappa_s = 0\) have no coupling axis and remain massless
  (S3).
\item
  \textbf{Electromagnetic scattering and pair annihilation}: The photon
  changes only the dynamic state \(\Delta\mathbf{k}\). The relative sign
  of the \(t\)-axis components determines attraction / repulsion. Pair
  annihilation of a direction-constrained wave and its anti-wave is the
  conversion \(\kappa = 2 \to \kappa = 0\) (S4).
\item
  \textbf{Gravity and wave packet deformation}: The graviton acts on
  both the \((t, R)\) axes and couples to all waves with
  \(\Delta\mathbf{k} \neq 0\). Under the conserved quantity, energy
  transfer deforms the amplitude and width of a wave packet inversely
  (S5).
\item
  \textbf{Strong force and confinement}: Mediates \(Q\)-axis transitions
  (color charge exchange), and the color neutrality condition provides
  the condition for forming an independent subjective space (S6).
\item
  \textbf{Weak force and CP violation}: Mediates \(t\)-axis displacement
  (flavor transition), and \(k_t \neq 0\) is the origin of CP violation
  and the left-handed chirality selection rule (S7).
\item
  \textbf{Wave packet deformation and causal structure}: Wave packet
  collapse is a conservation-law-governed deformation, and
  ``observation'' is unnecessary. Interactions mediated by spherical
  waves with \(k_t = 0\) are causally symmetric, and the path,
  localization, and variational principle are retroactively constructed
  (S8).
\end{enumerate}

The correspondence with Table A of {[}4{]} is shown below:

{\def\LTcaptype{none} % do not increment counter
\begin{longtable}[]{@{}
  >{\raggedright\arraybackslash}p{(\linewidth - 10\tabcolsep) * \real{0.1607}}
  >{\raggedright\arraybackslash}p{(\linewidth - 10\tabcolsep) * \real{0.1429}}
  >{\raggedright\arraybackslash}p{(\linewidth - 10\tabcolsep) * \real{0.1607}}
  >{\raggedright\arraybackslash}p{(\linewidth - 10\tabcolsep) * \real{0.3036}}
  >{\raggedright\arraybackslash}p{(\linewidth - 10\tabcolsep) * \real{0.1250}}
  >{\raggedright\arraybackslash}p{(\linewidth - 10\tabcolsep) * \real{0.1071}}@{}}
\toprule\noalign{}
\begin{minipage}[b]{\linewidth}\raggedright
Interaction
\end{minipage} & \begin{minipage}[b]{\linewidth}\raggedright
Spherical wave
\end{minipage} & \begin{minipage}[b]{\linewidth}\raggedright
Non-zero axis
\end{minipage} & \begin{minipage}[b]{\linewidth}\raggedright
Effect on direction-constrained wave
\end{minipage} & \begin{minipage}[b]{\linewidth}\raggedright
\(k_t\)
\end{minipage} & \begin{minipage}[b]{\linewidth}\raggedright
Causality
\end{minipage} \\
\midrule\noalign{}
\endhead
\bottomrule\noalign{}
\endlastfoot
Higgs coupling & Standing wave \(H^0\) & Spatial 1-axis & Decrease of
\(v\) (mass) & 0 & -- \\
Electromagnetic & \(\gamma\) & \(R\) & \(\Delta\mathbf{k}\) change
(scattering) & 0 & Retroactive \\
Strong & \(g\) & \(Q\) & \(Q\) change (color exchange) & 0 &
Retroactive \\
Weak & \(W^{\pm}\) & \(t\) & \(k_t\) change (flavor transition) &
\(\neq 0\) & Causal \\
Gravity & \(G\) & \(t, R\) & \(\Delta\mathbf{k}\) change (acceleration)
& \(\neq 0\) & Causal \\
\end{longtable}
}

All claims of this paper are stated as geometric facts, and the physical
correspondences are presented as ``connectable interpretations.''

\begin{center}\rule{0.5\linewidth}{0.5pt}\end{center}

\subsection{References}\label{references}

{[}1{]} Noriaki Kihara, ``Shape-Invariant Standing Waves on Closed
Spheres: Stability by Dimension and the Geometric Necessity of 3+1
Spacetime,'' 2026.

{[}2{]} Noriaki Kihara, ``Shape-Invariant Traveling Waves on a
4-Dimensional Integer Lattice: Existence Conditions, Conservation
Structure, and Self-Consistent Loop,'' 2026.

{[}3{]} Noriaki Kihara, ``Four Modes of Shape-Invariant Waves: Spin,
Chirality, and Statistics Determined by Wave Vector Structure,'' 2026.

{[}4{]} Noriaki Kihara, ``Set Structure of a 6-Dimensional
Hyperrectangle and Its Combinatorial Properties,'' Zenodo, 2025. DOI:
10.5281/zenodo.19688521.

{[}5{]} Noriaki Kihara, ``Three-Layer Model of the Universe via Central
Projection,'' Zenodo, 2025. DOI: 10.5281/zenodo.15083729.
\end{document}
